%% GENERATED from PAPER_SUBMISSION.md (pandoc) -- hand-fixed conversion; see BUILD_NOTES.md
\section{Introduction}\label{sec:1}

Sequential recommendation models predict a user's next item from their interaction
history. SASRec \citep{kang2018sasrec} and BERT4Rec \citep{sun2019bert4rec}
established Transformer-based baselines. Later systems add language-derived item
representations \citep{ding2021zesrec,hou2022unisrec,li2023recformer,hou2024blair}
or discrete semantic IDs \citep{rajput2023tiger,yang2024liger}.

Amazon Reviews 2023 results span incompatible preprocessing variants. In particular,
the 0-core reference pipeline and the 5-core protocol used by recent generative
retrievers do not define interchangeable benchmarks. Our experiments use the latter:
fixed 5-core splits, leave-last-two-out histories, and full-catalog ranking.

Recommender evaluation is vulnerable to mistuned baselines, incompatible protocols,
and irreproducible result selection
\citep{ferraridacrema2019worrying,ferraridacrema2021troubling}. We therefore pair
the method with a per-paper discipline: git-frozen protocols, a fail-closed graph that
recomputes 195 reported cells across 19 claim families, environment-caveated comparator
regeneration, and symmetric adjudication. Failures are retained. In particular, the
original Office\_Products protocol remains VOID after its own floor check exposed an
environment mismatch, although a redesigned fresh-seed protocol later passed
(\S\ref{sec:5.2}; Appendix A.0).

Using an HSTU-style pure-PyTorch encoder based on \citet{zhai2024hstu}, we ask whether
one small causal filtering component survives matched controls and can be reused without
category-specific tuning. Our contribution ledger is deliberately limited to three
items:

\begin{enumerate}
\item \textbf{A canonical causal FIR module.} We adapt earlier bidirectional frequency
filters to a strictly left-causal, depthwise residual that is identity-initialized and
gradient-active under the all-position objective (\S\ref{sec:3}). This is a modular
component, not a new recommender architecture.
\item \textbf{Matched internal evidence and a mechanism boundary.} Learned
FIR--identity estimates are positive on Musical\_Instruments (+0.002265, 95\% CI
[0.001928, 0.002602]) and under one untuned configuration on
Industrial\_and\_Scientific (+0.002110 [0.001820, 0.002399]) and
CDs\_and\_Vinyl (+0.006150 [0.005849, 0.006450]) (\S\ref{sec:5.2}). Because
these studies reuse outcome-visible splits, they are internal robustness evidence, not
independent confirmation. On Musical\_Instruments, learned taps do not separate from a
shared causal filter, but they beat a parameter-matched current-position-only placebo by
+0.001941 [+0.001788, +0.002095], while the placebo--identity interval spans zero. This
discriminates learned FIR from that compound placebo; it does not isolate temporal access,
establish unique necessity of per-channel taps, show cross-domain generalization, or provide independent confirmation. A frequency-5-heavy Musical\_Instruments
text-tail effect is secondary and does not establish cross-dataset heterogeneity
(\S\ref{sec:5.3}--\S\ref{sec:5.4}).
\item \textbf{An auditable evaluation record.} Released protocols, artifacts,
sidecars where available, comparator-regeneration code, and a strict rebuild preserve
positive, null, deviated, and VOID outcomes under the same reporting rule
(\S\ref{sec:8}). The apparatus certifies reconstruction and protocol history; it does
not strengthen the estimand.
\end{enumerate}

TAPE, late fusion, sparse-warm redistribution, titrations, probe screens, and earlier
Beauty results are supporting studies. All borrowed architectures and training recipes
are attributed in \S\ref{sec:2}.
